\begin{textsample}{POS Dim 1 – vem\_ed\_30 – Score 54.00 – vem\_ed\_30/t156.txt}  \label{ex:f1_pos_026}
Osvaldo Succi Junior na França: International Digital Week Osvaldo Succi Junior, coordenador da área de Apoio à Internacionalização do Ensino Superior da Divisão de Extensão e Pesquisa (Depes) da Coordenadoria Geral de Ensino Superior de Graduação (CGESG) do CPS, \textbf{foi} convidado para a International Digital Week 2025.
O evento, organizado pela Université Côte D´Azur, ocorreu de 16 a 20 de junho na universidade em Nice, França.
O tema do ano \textbf{foi} Digital Transformation in Higher Education (Transformação Digital na Educação Superior): https://univ-cotedazur.eu/events/digital-week.
Eventos como esse não apenas \textbf{promovem} a \textbf{troca} de \textbf{conhecimentos} técnicos, mas também: Fortalecem redes \textbf{colaborativas} entre instituições internacionais; Democratizam o acesso a práticas educacionais inovadoras; Aceleram a adoção de tecnologias emergentes no ensino; Preparam educadores para os desafios da \textbf{era} digital.
Em 16 de junho, Succi Junior apresentou um workshop sobre a revolução da Inteligência Artificial nos projetos COIL (Collaborative Online International Learning), \textbf{que} nas Fatecs do CPS recebem o nome de Projetos Colaborativos Internacionais (PCIs).
A \textbf{proposta} foi discutir o \textbf{papel} da IA nesses projetos e seu potencial para dinamizar experiências interculturais e \textbf{interdisciplinares} na formação \textbf{profissional} – tema abordado no Artigo de \textbf{Opinião} deste \textbf{número} de VEm. \textbf{continuação} No dia 17 de junho, a professora Luciane Stallivieri (UFSC) conduziu uma instigante discussão sobre o conceito de internacionalização responsável no workshop"Implementing Successful COIL and Virtual Exchanges VE / COIL Roadmap", ("Implementando um Roteiro para COIL e Intercâmbios Virtuais \textbf{bem-sucedidos}"), enfatizando a \textbf{importância} de abordagens éticas e \textbf{sustentáveis}.
Complementando essa fundamentação teórica, Succi Junior, com sua experiência na \textbf{coordenação} de Projetos Colaborativos Internacionais (PCIs), auxiliou os participantes a explorar a aplicação \textbf{prática} de alguns desses \textbf{princípios}, guiando-os na elaboração de \textbf{estratégias} concretas para o \textbf{desenvolvimento} de \textbf{intercâmbios} virtuais \textbf{bem-sucedidos}.
Em 19 de junho, apresentou um pôster sobre o PCI conduzido no primeiro semestre de 2025 pelos professores Leandro Carvalho e Cristiane Pellizzon (Fatec Itatiba) com o professor Filip Burgelman, da Thomas More University of Applied Sciences (Bélgica).
Nesse projeto, os estudantes desenvolveram um game \textbf{educativo} para conscientizar crianças de 6 a 9 anos sobre a preservação dos recursos hídricos.
O trabalho \textbf{foi} eleito o “Melhor Pôster do Dia” pelos participantes do evento. “Esse prêmio eleva o perfil do Centro Paula Souza como instituição capaz de \textbf{produzir} resultados de excelência em \textbf{contextos} globais”, comenta Succi Junior.
O PCI"Learning Games \textbf{for} Environmental Awareness: Educating Children on Clean Water and Nature Preservation"("Jogos \textbf{educativos} para a conscientização \textbf{ambiental}: educando as crianças sobre a água potável e a preservação da natureza") \textbf{é} um excelente exemplo de alinhamento à missão da Depes.
Os professores conseguiram instigar os alunos a desenvolver jogos digitais para crianças, a \textbf{partir} de contatos articulados pela Fatec Itatiba com as ONGs JAPPA e SIBES e pela Thomas More University com a ONG City to Ocean.
Essa \textbf{conexão} entre instituições acadêmicas e organizações da sociedade civil \textbf{representa} um exemplo \textbf{educativo} \textbf{que} vai além da sala de aula, \textbf{permitindo} \textbf{que} os estudantes compreendam o \textbf{impacto} real de seu trabalho na comunidade e desenvolvam uma consciência social mais \textbf{ampla}. \textbf{continuação} O projeto ilustra perfeitamente como os PCIs \textbf{promovem} a formação integral dos estudantes.
Ao trabalhar em colaboração com parceiros internacionais, os alunos desenvolvem não apenas \textbf{competências} técnicas \textbf{específicas} de suas áreas, mas também \textbf{habilidades} interculturais, \textbf{capacidade} de comunicação em \textbf{contextos} globais, pensamento \textbf{crítico} sobre \textbf{questões} \textbf{ambientais} universais e sensibilidade para problemáticas sociais \textbf{que} transcendem fronteiras nacionais.
Os PCIs das Fatecs, coordenados pela área de Apoio à Internacionalização do Ensino Superior da Depes, refletem o esforço em \textbf{promover} educação para a cidadania global aos fatecanos, \textbf{seguindo} as diretrizes da Assessoria de Relações Internacionais do CPS, com um \textbf{foco} \textbf{específico} nas \textbf{ações} de Internacionalização em Casa. “Os PCIs têm uma função primordial de \textbf{promover} a qualidade e a \textbf{capacidade} dos alunos das Fatecs em palcos internacionais, preparando-os não apenas para o \textbf{mercado} de trabalho \textbf{local}, mas para uma cidadania global \textbf{ativa} e consciente”, ressalta Succi Junior. \textbf{continuação} Em 20 de junho, Succi Junior participou da mesa-redonda intitulada “Entrepreneurship \& AI: Redesigning the Future of Higher Education” (“Empreendedorismo \& IA: redesenhando o \textbf{futuro} da Educação Superior”), com o especialista em educação online Simon Carolan (França) e Robin Couture-Matte, professor da Téluq University (Canadá).
A mediação \textbf{foi} de Julia Reinhard, da Haaga-Helia University of Applied Sciences (Finlândia).
Nesse mesmo dia, falou sobre “COIL as a Launchpad for Entrepreneurial Collaboration” (“COIL como plataforma para colaboração empreendedora”), aprofundando as possibilidades dos projetos COIL para fomentar iniciativas de empreendedorismo internacional entre estudantes e instituições parceiras. “A Digital Week 2025 \textbf{foi} um marco importante: \textbf{permitiu} reforçar laços com a comunidade acadêmica global, explorar o potencial disruptivo da digitalização e da inteligência \textbf{artificial} na educação \textbf{superior} e, acima de tudo, posicionar o Centro Paula Souza como ator \textbf{relevante} na vanguarda da Internacionalização em Casa”, resume.

% matched lemmas: ambiental, amplo, artificial, ativo, ação, bem-sucedido, capacidade, colaborativo, competência, conexão, conhecimento, contexto, continuação, coordenação, crítico, desenvolvimento, educativo, específico, estratégia, foco, futuro, habilidade, impacto, importância, intercâmbio, interdisciplinar, local, mercado, número, opinião, papel, partir, permitir, princípio, produzir, profissional, promover, proposta, prático, que, questão, relevante, representar, seguir, ser, superior, sustentável, troca
\end{textsample}
